% Transcription — fonds Alexandre Grothendieck, shelfmark 81, batch 2 (pages 21-40).
% Established from the facsimile archives/batches/81/batch-02.pdf (a mirror of
% https://grothendieck.umontpellier.fr/81.pdf), per the skill
% transcribe-grothendieck. The batch file carries no cover sheet: PDF page N
% is archivists' page 20+N; the Montpellier footer prints « 21 » ... « 40 »,
% and the pencilled numbers bottom-left agree wherever legible (« 31 »,
% « 32 », « 33 », « 34 », « 35 », « 36 », « 37 », « 38 », « 39 » among
% others). He did not paginate these leaves; his own formula numbers (1) to
% (27) run through pages 31-38.
%
% Pass: Opus 5.5 (claude-opus-5-5), 2026-09-23 — first pass, unchecked against the pages by a human.
%
% Eighteen of the twenty pages carry a \page. Absent: page 30, a blank
% (pink) sheet; page 40, an otherwise blank sheet inscribed in pencil, top
% right, in his hand, « relations d'équi. non parallèles » — a cover for the
% run that opens batch 3 (page 41 ff.; « Relation d'équivalence non
% parallèle » is the heading of its pages 48-50). By the rule for covers its
% words head the next section, which lies in batch-03.fr.tex and opens on
% them.
%
% What the batch is. Three runs, all on non parallel « binary
% decompositions » of a space X into two closed « hemispheres » meeting in
% an « equator » — the combinatorial frame for cutting a disc along
% disjoint arcs.
%   21-24  End of a run begun in batch 1: a family (E_i)_{i∈I} of disjoint
%          « equators », each cutting X into Σ_i^+ and Σ_i^-. Page 21: for
%          V = X \ ∪_{α∈K} Σ̂_α^+ (the Σ̂_α^+ disjoint), V̄ ∩ E_α =
%          closure(Σ_α^-) ∩ E_α ≠ ∅ (a boxed, cancelled attempt, then a
%          two-line proof). Page 22: V̄ ∩ E_β = ∅ for β ∉ K; Lemma (each E_i
%          lies in one Σ_α^+); Corollary: the partition Π_K^0 consists of the
%          Σ_α^+ and V. Page 23: the classes U_s (s ∈ S) of Π_I^0, the
%          relation R ⊂ S × I (closure(U_s) ∩ E_i ≠ ∅); Lemma: (S, I, R) is a
%          tree. Page 24 stops after a third of a page.
%   25-29  A fresh start: « disposition binaire » of a set, hemispheres,
%          equator; two dispositions « non parallèles »; the chain
%          Σ_1^+ ⊃ E_1 → Σ_12 ← E_2 → Σ_2^+ and a long oblique margin on the
%          degenerate case; page 27: from a tree Γ = (S, A, R) and families
%          (U_s), (E_a), X = ∐ U_s ∐ ∐ E_a with one binary decomposition per
%          edge; conditions (*) end-vertices, (**) and (***) order-2
%          vertices; page 28: its picture as a graph with a point on each
%          edge (barycentric subdivision); page 29: the traces on S* = S ∐ A'.
%   31-39  « Version II », titled « Décompositions n-aires d'un espace
%          topologique », numbered 1.1-1.3 with formulas (1)-(27): binary
%          decomposition (closed hemispheres, equator, Σ^{0ε}), 1.2 non
%          parallel pairs, the inclusion diagram (14) with its three
%          cartesian-cocartesian squares, (16) Σ_12 ≠ ∅, Proposition
%          (uniqueness of the disjoint pair), the cases (20) where the four
%          sets are not distinct, Déf « strictement non parallèles » (21), the
%          poset (22) on four elements; 1.3 families (Π_α)_{α∈A} mutually
%          non parallel (« décomposition n-aire »), C = ∐ ω_α with its
%          fixed-point-free involution σ and order (26), the increasing map
%          (27) C → closed parts of X; Déf « quartette ordonné »; page 39:
%          Proposition, C_{≥x} and C_{≤x} totally ordered — proof broken off
%          mid-formula, the rest of the page blank.
%
% Notation fixed for the batch. \widehat{\Sigma} for his Σ with a hat
% (pages 21-24, the closed hemisphere), \Sigma^{0\varepsilon} for his
% « Σ^{0ε} » (hemisphere minus equator), \mathfrak{P}_2 for his script P_2,
% \smallsetminus for set difference. The run 21-29 is a fast hand whose
% formulas carry and whose prose often does not; 31-39 are written fair,
% with oblique marginal notes that are largely uncertain. His tree diagram
% (14) on page 33 and the four-point order on page 37 are commutative-diagram
% like and go into tikzcd; the pictures (a segment with three points, a cross
% of arrows) are described in \note{}.
%
% The dating is the inventory's own for the folder, copied verbatim. Its
% group, [69-89] « Géométrie et topologie combinatoire ».
\documentclass[11pt,a4paper]{article}
% Le préambule commun à toutes les transcriptions du fonds Grothendieck.
%
% Il tient en une idée : **l'incertitude se compose, elle ne se résout pas.**
% Une transcription qui lisse un mot douteux a détruit la seule chose pour
% laquelle on la faisait — un lecteur doit pouvoir savoir, à chaque endroit,
% ce qui était sur le feuillet et ce qui a été ajouté. Les six macros qui
% suivent sont donc l'appareil critique, et non de la décoration.
%
% Les mêmes commandes sont reconnues par `scripts/render.mjs`, qui produit la
% vue de lecture du site. Une macro ajoutée ici doit l'être aussi là-bas,
% faute de quoi le rendu échouera bruyamment — ce qui est voulu : un rendu
% silencieusement faux est pire qu'un rendu absent.

\ProvidesPackage{grothendieck}[2026/08/08 Transcriptions du fonds Grothendieck]

\usepackage{amsmath,amssymb,amsthm}
\usepackage{mathrsfs}
\usepackage{stmaryrd}
% Les diagrammes commutatifs sont partout dans ce fonds : c'est la langue
% même de la géométrie algébrique de Grothendieck, et une transcription qui
% les rendrait en prose ne transcrirait rien.
\usepackage{tikz-cd}
% Français par défaut — c'est la langue des feuillets — mais l'anglais est
% chargé aussi : la traduction et le résumé sont des documents anglais, et la
% typographie française (espaces avant les deux-points) y serait une faute.
% Ces documents basculent par \selectlanguage{english} en tête de corps.
\usepackage[english,french]{babel}
\usepackage{fontspec}
\usepackage{geometry}
\geometry{a4paper,margin=25mm,marginparwidth=28mm}
\usepackage{marginnote}
\usepackage{xcolor}
% ulem, not soul. `soul`'s \st and \ul are built on a character-by-character
% scan that XeLaTeX's font machinery breaks: the document fails with an
% undefined control sequence rather than degrading. ulem does the same two jobs
% and is Unicode-engine safe. `normalem` stops it hijacking \emph, which the
% transcriptions use for Grothendieck's own emphasis.
\usepackage[normalem]{ulem}
\usepackage{hyperref}
\hypersetup{colorlinks=true,linkcolor=black,urlcolor=black,pdfborder={0 0 0}}

% --- Métadonnées du lot -----------------------------------------------------
% Reprises telles quelles par le script de rendu, qui les affiche en tête de
% la vue de lecture. Elles fixent ce que le fichier prétend couvrir : sans
% elles, une transcription flotte sans point d'attache dans un fonds de seize
% mille feuillets.

\newcommand{\folder}[1]{\gdef\grothFolder{#1}}
\newcommand{\batch}[1]{\gdef\grothBatch{#1}}
\newcommand{\leaves}[2]{\gdef\grothFirst{#1}\gdef\grothLast{#2}}
\newcommand{\foldertitle}[1]{\gdef\grothFoldertitle{#1}}
\gdef\grothFolder{?}\gdef\grothBatch{?}\gdef\grothFirst{?}\gdef\grothLast{?}\gdef\grothFoldertitle{}

% --- Le résumé --------------------------------------------------------------
%
% Il ouvre la lecture modernisée, sous le titre « Résumé », et il est composé
% en petit. Ce n'est pas un ornement : c'est le seul passage du document qui
% ne soit pas des mathématiques, et un lecteur doit pouvoir le reconnaître
% avant de l'avoir lu — puis le sauter s'il connaît déjà le sujet, ou s'y
% arrêter s'il ne le connaît pas. Le filet qui le suit marque la reprise du
% corps.

\newenvironment{resume}
  {\small\setlength{\parskip}{0.45em}}
  {\par\medskip\hrule\medskip}

% --- Mots-clés ---------------------------------------------------------------
%
% En anglais, en clôture du résumé de la lecture modernisée : le vocabulaire
% moderne sous lequel chercher ce que les feuillets construisent. Le manifeste
% du site (`npm run manifest`) les extrait pour étiqueter la cote entière —
% cette ligne est la source unique des tags, il n'y a pas de fichier à côté.
\newcommand{\keywords}[1]{\par\smallskip\noindent
  {\footnotesize\textsc{Keywords} --- \textit{#1}}\par}

% --- L'appareil critique ----------------------------------------------------

% Le numéro de feuillet, en marge. C'est l'ancre qui permet de régler un
% désaccord sur une formule en regardant *un* feuillet plutôt qu'une chemise
% de six cents. Le site s'en sert aussi pour tourner les pages du fac-similé
% quand on fait défiler la transcription.
\newcommand{\leaf}[1]{%
  \marginnote{\footnotesize\color{gray!70}\textsf{#1}}%
  \label{leaf:#1}\ignorespaces}

% Illisible : on ne devine pas. Un mot inventé qui se lit comme les autres est
% le pire résultat possible.
\newcommand{\ill}{\textcolor{red!60!black}{[\,\ldots\,]}}

% Lecture douteuse : le mot est donné, et signalé comme incertain.
\newcommand{\uncertain}[1]{\uline{#1}}

% Ajout de l'éditeur — une lettre manquante, un mot sous-entendu.
\newcommand{\add}[1]{\textcolor{blue!50!black}{[#1]}}

% Biffé par Grothendieck lui-même. Ce qu'il a rayé fait partie du document :
% une rature dit souvent où la pensée a changé de direction.
\newcommand{\struck}[1]{\sout{#1}}

% Note du transcripteur, sur l'état matériel du feuillet.
\newcommand{\note}[1]{\par\smallskip{\small\color{gray!60!black}\textit{[#1]}}\par\smallskip}

% Note marginale de Grothendieck, distincte du corps du texte.
\newcommand{\marginal}[1]{\par\smallskip\noindent\hspace{1em}%
  {\small\color{gray!60!black}#1}\par\smallskip}

% --- Titre ------------------------------------------------------------------

\AtBeginDocument{%
  \begin{center}
    {\large\bfseries Fonds Alexandre Grothendieck --- cote n\textsuperscript{o}~\grothFolder}\\[2pt]
    {\itshape\grothFoldertitle}\\[4pt]
    {\small feuillets \grothFirst--\grothLast\ (lot \grothBatch)}\\[2pt]
    {\footnotesize\color{gray!60!black}Transcription établie d'après le fac-similé numérisé
     par l'Université de Montpellier.\\
     Les lectures incertaines sont soulignées, les illisibles notées [\,\ldots\,],
     les ajouts entre crochets.}
  \end{center}
  \vspace{1em}}

\endinput


\folder{81}
\batch{2}
\pages{21}{40}
\dating{[à partir de 1981]}
\watermark{Édition de démonstration}
\foldertitle{Immersions du disque et de la circonférence : notes manuscrites (s.d.).}

\begin{document}

\page{21}
\struck{Un procédé par récurrence sur $\operatorname{card} K$, pour prouver
que} On note d'abord que (si $V=X\smallsetminus\bigcup_{\alpha\in
K}\widehat{\Sigma}_{\alpha}^{+}$, les $\widehat{\Sigma}_{\alpha}^{+}$
($\alpha\in K$) étant disjoints)
\note{sous l'accolade qui souligne $X\smallsetminus\bigcup_{\alpha\in
K}\widehat{\Sigma}_{\alpha}^{+}$, il écrit « $=\bigcap_{\alpha\in
K}\Sigma_{\alpha}^{-}$ »}
\[
  \left\{
  \begin{array}{l}
    \overline{V}\cap E_{\alpha}\neq\emptyset\\
    V\cap\Sigma_{\alpha}^{-}\neq\emptyset\quad\text{dès que}\quad
    V\cap\Sigma_{\alpha}^{+}=\emptyset
  \end{array}
  \right.
\]
\marginal{trivial}
\note{« trivial » est écrit à gauche de la seconde ligne, dont le $\neq$
remplace un $=$ biffé}

De façon précise, je dis que
\[
  \overline{V}\cap E_{\alpha}=\overline{\Sigma_{\alpha}^{-}}\cap E_{\alpha}
  \qquad(\text{qui est donc}\neq\emptyset)
\]
\note{au-dessous, une ligne biffée d'un trait :
« $=\overline{V}\cap\widehat{\Sigma}_{\alpha}^{+}$ » sous le membre de gauche,
« $\overline{\Sigma_{\alpha}^{-}}\cap\widehat{\Sigma}_{\alpha}^{+}$ » sous
celui de droite}

\struck{Comme $V\subset\Sigma_{\alpha}^{-}$, on a $V\cap
E_{\alpha}\subset\overline{\Sigma_{\alpha}^{-}}\cap E_{\alpha}$, prouvons
l'inclusion inverse}
\struck{$\overline{\Sigma_{\alpha}^{-}}\cap E_{\alpha}\subset
\overline{V}\cap E_{\alpha}$.}
\struck{Soit $x\in\overline{\Sigma_{\alpha}^{-}}\cap E_{\alpha}$, prouvons
$x\in\overline{V}\cap E_{\alpha}$ i.e.\ $x\in\overline{V}$, soit $W$ un
voisinage de $x$, montrons qu'il rencontre \ill{}.}
\struck{Or \ill{} par construction cette égalité comme égalité dans
l'espace}
\struck{$\widehat{\Sigma}_{\alpha}^{-}=\Sigma_{\alpha}^{-}\cup E_{\alpha}$}
\struck{dans laquelle on a deux parties fermées disjointes}
\note{tout ce passage est encadré d'un trait et barré de plusieurs longues
diagonales ; devant « $=\widehat{\Sigma}_{\alpha}^{-}$ » un $X$ est biffé,
et « égalité » remplace un mot biffé en début de ligne}

Or dans $X$ on a les deux parties fermées \emph{disjointes}
\[
  Y=\widehat{\Sigma}_{\alpha}^{+},\qquad
  Z=\bigcup_{\beta\in K\smallsetminus\{\alpha\}}\widehat{\Sigma}_{\beta}^{+}
\]
et on doit prouver
\[
  \overline{\bigl(X\smallsetminus(Y\cup Z)\bigr)}\cap Y
  =\overline{(X\smallsetminus Y)}\cap Y
\]
ce qui est immédiat.

\page{22}
Soit maintenant $\beta\in I\smallsetminus K$, prouvons que l'on a
\[
  \overline{V}\cap E_{\beta}=\emptyset .
\]
En effet, $V$ \struck{contient} ne rencontre aucun $E_{i}$ ($i\in I$) donc
\[
  E_{\beta}\subset X\smallsetminus V=\bigcup_{\alpha\in
  K}\widehat{\Sigma}_{\alpha}^{+},
\]
\struck{donc si \ill{}} je dis que si $\beta\in I\smallsetminus K$, on a en fait
\[
  E_{\beta}\subset\bigcup\Sigma_{\alpha}^{+}\qquad(\text{ouvert disjoint de
  }V)
\]
qui prouve l'assertion — mieux, je dis que $\exists\,\alpha\in K$ tel que
$E_{\beta}\subset\Sigma_{\alpha}^{+}$. Cela résulte du
\marginal{c'est trivial, du fait que les $E_{i}$ ($i\in I$) sont disjoints}

\emph{Lemme} Soit $K\subset I$, \struck{tel que les} et des
$\Sigma_{\alpha}^{+}$ ($\alpha\in K$) tels que
\[
  \alpha\neq\beta\in K,\ \alpha\neq\beta\ \Rightarrow\ \Sigma_{\alpha}^{+}\cap\Sigma_{\beta}^{+}=\emptyset .
\]
Soit $i\in I\smallsetminus K$ tel que $E_{i}\subset\bigcup\Sigma_{\alpha}^{+}$,
alors \struck{$E_{i}$ est contenu dans un} $\exists\,\alpha\in K$, tel que
$E_{i}\subset\Sigma_{\alpha}^{+}$.

Cela \struck{signifie que} \uncertain{provient} \ill{}
\note{la phrase reste inachevée}

\emph{Corollaire} Avec les hyp.\ précédentes sur
$(\Sigma_{i}^{\alpha})_{\alpha\in K}$, la partition $\Pi_{K}^{0}$ est formée
des $\Sigma_{\alpha}^{+}$ ($\alpha\in K$) et de
\[
  V=X\smallsetminus\bigcup_{\alpha\in K}\widehat{\Sigma}_{\alpha}^{+} .
\]
\note{la réunion est écrite avec le signe de réunion disjointe}

Immédiat par récurrence sur $\operatorname{card} K\geqslant0$.

\page{23}
Considérons la partition $\Pi_{I}^{0}$ de $X\smallsetminus\bigcup_{i\in
I}E_{i}$, soit $S$ l'ens.\ des classes, notées $U_{s}$ ($s\in S$).
\struck{On} Considérons la \struck{relation $R$ entre} \struck{$\forall
i\in I$,} \add{relation $R$ entre} $S$ et $I$
\[
  (s,i)\in R\ \overset{\text{déf}}{\Longleftrightarrow}\ \overline{U_{s}}\cap E_{i}\neq\emptyset .
\]

\emph{Lemme} Cette relation fait de $(S,I,R)$ un graphe strict ayant $S$
comme ens.\ de sommets, $I$ comme ens.\ d'arêtes. \struck{Cette relation}
Ce graphe est un \emph{arbre}.

Le fait que c'est un graphe strict signifie que a) $\forall i\in I$,
\struck{il y a} il y a exactement deux $U_{s}$ tels que
$\overline{U_{s}}\cap E_{i}\neq\emptyset$, et b) $i$ est caractérisé par
l'une de ces deux $U_{s}$.

Soient $\Sigma_{i}^{+}$, $\Sigma_{i}^{-}$ les deux comp.\ conn.\ de
$X\smallsetminus E_{i}$, alors a) se précise ainsi : il y a exactement
\struck{une} une $U_{s_{+}}$ contenue dans $\Sigma_{i}^{+}$ telle que
$\overline{U_{s_{+}}}\cap E_{i}\neq\emptyset$ (et de même pour
$\Sigma_{i}^{-}$, \ill{} définit une $U_{s_{-}}\subset\Sigma_{i}^{-}$ telle
que $\overline{U_{s_{-}}}\cap E_{i}\neq\emptyset$),

\page{24}
\struck{Soit en effet} Si on a $U_{s}\subset\Sigma_{i}^{-}$,
$\overline{U_{s}}\cap E_{i}\neq\emptyset$, \struck{ainsi} alors on a
\[
  U_{s}=X\smallsetminus\bigcup_{\alpha\in K}\widehat{\Sigma}_{\alpha}^{+}
  \qquad\text{où } i\in K\subset I
\]
$\forall\alpha\in K$
$\Sigma_{\alpha}^{+}\in\pi_{0}(X\smallsetminus E_{\alpha})$ \struck{tel},
et les $\Sigma_{\alpha}^{+}$ mut\add{uellemen}t disjoints.
\note{dans $U_{s}\subset\Sigma_{i}^{-}$, le $\Sigma$ est repassé et son
exposant mal formé ; la réunion est écrite avec le signe de réunion
disjointe}

Il faut voir que \struck{$K$} $K\smallsetminus\{i\}$ et les
$\Sigma_{\alpha}^{+}$ ($\alpha\in K\smallsetminus\{i\}$) sont déterminés en
termes de $i$ et de $\Sigma_{i}^{+}$ i.e.\ $\Sigma_{i}^{-}$. Je dis que
\[
  K\smallsetminus\{i\}=\{\alpha\in I\mid E_{\alpha}\subset\Sigma_{\alpha}^{-}\}
\]
\note{l'indice du dernier $\Sigma$ est bien un $\alpha$ sur la page ; on
attendrait $\Sigma_{i}^{-}$. La page s'arrête là, les deux tiers inférieurs
restent blancs}

\page{25}
Soit $X$ un ens. Une \emph{disposition binaire} de $X$ est un ens.\ de deux
parties de $X$
\[
  (\Sigma^{\alpha})_{\alpha\in\omega}\qquad(\text{les “hémisphères”})
\]
\note{l'indice de la famille est mal formé ; on lit $\alpha\in\omega$ ou
$\alpha\in\varepsilon$}
dont la réunion est $X$, et dont \struck{l'intersection} aucune n'est
contenue dans l'autre — on désigne par $E$ (“équateur”) l'intersection (on
n'exclut pas qu'elle soit vide). Quand $X$ est un espace topologique, on
\uncertain{suppose} \ill{} \ill{} que \struck{hémisphères} les
$\Sigma_{\alpha}$ sont fermés, (et \add{de} plus, le plus souvent,
\uncertain{connexes} \ill{} et \uncertain{si} $E\neq\emptyset$) ; si les
$\Sigma^{\alpha}$ connexes et $E\neq\emptyset$, $X$ est connexe,
\struck{\ill{}} si \struck{les $\Sigma^{\alpha,0}$ connexes}, on conclut
\struck{aussitôt} que $\overline{\Sigma^{\alpha,0}}\cap E\neq\emptyset$ (car
$\overline{\Sigma^{\alpha,0}}\cap E=\overline{\Sigma^{\alpha,0}}
\smallsetminus\Sigma^{\alpha,0}$, et \uncertain{si} $\Sigma^{\alpha,0}$ une
partie à la fois ouverte et fermée $\neq\emptyset$ et $\neq X$)
\note{interligne, au-dessus de « le plus souvent » :
« $\Sigma^{\alpha}$ et $\Sigma^{\alpha}\smallsetminus E=\Sigma^{\alpha,0}$ »
; la fin de la phrase, depuis « connexes », est entourée d'un trait qui la
rattache à cette insertion}

\uncertain{car} $\Sigma^{\alpha,0}$ est un \uncertain{ouvert} \ill{}
\note{début de ligne laissé en suspens}

\marginal{diagramme \struck{\ill{}} : un segment,
$\Sigma^{-}\supset E\subset\Sigma^{+}$}
\note{dans la marge gauche, en oblique : un segment entre deux points, avec
dessous « $\Sigma^{-}\supset E\subset\Sigma^{+}$ »}

Considérons deux dispositions binaires $\{\Sigma_{1}^{+},\Sigma_{1}^{-}\}$,
$\{\Sigma_{2}^{+},\Sigma_{2}^{-}\}$ \struck{distinctes} de $X$, on dit
qu'elles sont \emph{non parallèles} \add{si elles sont \emph{distinctes},
et} si on peut trouver $\Sigma_{1}^{+}$ dans l'une, $\Sigma_{2}^{+}$ dans
l'autre, tels que
\[
  \Sigma_{1}^{+}\cap\Sigma_{2}^{+}=\emptyset\qquad(\text{d'où }E_{1}\cap
  E_{2}=\emptyset)
\]
on a donc
\[
  \begin{array}{ll}
    \Sigma_{1}^{+}\subset\Sigma_{2}^{-} &
    \text{i.e.}\ \Sigma_{1}^{+}\cap\Sigma_{2}^{-}=\Sigma_{1}^{+}\\
    \Sigma_{2}^{+}\subset\Sigma_{1}^{-} &
    \text{i.e.}\ \Sigma_{2}^{+}\cap\Sigma_{1}^{-}=\Sigma_{2}^{+}
  \end{array}
\]
\[
  \Sigma_{12}\overset{\text{déf}}{=}\Sigma_{1}^{-}\cap\Sigma_{2}^{-}\supset
  E_{1}\cup E_{2},\qquad\Sigma_{1}^{-}\cap\Sigma_{2}^{-}\neq\emptyset
\]
\note{après $\Sigma_{1}^{+}\subset\Sigma_{2}^{-}$, un « $\neq\emptyset$ »
biffé}

\page{26}
cette dernière relation $\Sigma_{1}^{-}\cap\Sigma_{2}^{-}\neq\emptyset$ est
vraie \uncertain{même} si $E_{1}=E_{2}=\emptyset$, car \struck{\ill{}} si
\uncertain{on avait} de plus $\Sigma_{1}^{-}\cap\Sigma_{2}^{-}=\emptyset$, on
en conclurait que les deux \struck{part.} décompositions binaires
\struck{sont} (qui sont donc des partitions binaires) \uncertain{sont}
\emph{identiques}, ce que l'on \uncertain{exclut}. Cela montre que les
$\Sigma_{1}^{+}$, $\Sigma_{2}^{+}$ sont uniquement déterminés.

On a donc le diagramme d'inclusions
\[
  \Sigma_{1}^{+}\supset E_{1}\hookrightarrow\Sigma_{12}\hookleftarrow
  E_{2}\hookrightarrow\Sigma_{2}^{+}
\]
qu'on \uncertain{exprime} par un diagramme
\note{un segment, trois points marqués (les deux extrémités et le milieu),
et dessous
« $\Sigma_{1}^{+}\supset E_{1}\hookrightarrow\Sigma_{12}\hookleftarrow
E_{2}\hookrightarrow\Sigma_{2}^{+}$ »}
\marginal{NB $\left\{\begin{array}{l}
\Sigma_{1}^{-}=\Sigma_{12}\amalg_{E_{2}}\Sigma_{2}^{+}\\
\Sigma_{2}^{-}=\Sigma_{12}\amalg_{E_{1}}\Sigma_{1}^{+}
\end{array}\right.$}

Je me propose d'étudier les ensembles \struck{\ill{}}
$(\Pi_{\alpha})_{\alpha\in A}$ \struck{$\ill{}$, \ill{}} de décompositions
binaires de $X$ \struck{\ill{}} deux à deux non parallèles. \uncertain{Voici}
une façon \struck{très} générale \add{d'en} obtenir des \uncertain{exemples}.

\marginal{Ceci n'implique pas que les quatre ensembles
$\Sigma_{1}^{+},\Sigma_{1}^{-},\Sigma_{2}^{+},\Sigma_{2}^{-}$ soient
distincts, car on peut avoir (bien que $\Sigma_{1}^{+}\neq\Sigma_{2}^{+}$, ni
$\Sigma_{1}^{-}=\Sigma_{2}^{-}$, mais) $\Sigma_{1}^{+}=\Sigma_{2}^{-}$, qui
implique : $E_{1}=\Sigma_{12}$ et \uncertain{contraire} $E_{2}=\emptyset$
\ill{} (cela signifie, en termes de $(\Sigma_{1}^{+},\Sigma_{1}^{-})$, que
$\Sigma_{2}^{-}=\Sigma_{1}^{+}$,
$\Sigma_{2}^{+}=X\smallsetminus\Sigma_{1}^{+}=\Sigma_{1}^{-,0}$) \ill{}
symétrique. On \uncertain{demande} \ill{} \uncertain{en ce cas}, i.e.\
\uncertain{on} n'exige pas que les quatre ensembles
$\Sigma_{1}^{+},\Sigma_{1}^{-},\Sigma_{2}^{+},\Sigma_{2}^{-}$ soient
distincts \ill{} \ill{} les \struck{ens.} deux décompositions
\uncertain{sont strict.} non parallèles.}
\note{longue note de sa main, écrite en oblique dans toute la marge gauche,
d'une lecture très incertaine ; elle porte aussi un petit segment entre deux
points marqué « $\Sigma_{1}^{+}\supset E_{1}\subset\Sigma_{12}$ » et, sur une
ligne biffée, « $\Sigma^{+}_{i}\in\pi_{0}\ldots$ ». Vers la fin, la
correction « strict.\ » est ajoutée devant « non parallèles », ce qui
annonce la notion de la page 36}

\page{27}
Soit $\Gamma=(S,A,R)$ un arbre ayant $S$ comme ens.\ de sommets, $A$ comme
ens.\ d'arêtes, $R$ comme relation d'incidence [\uncertain{définie} par
$A\hookrightarrow\mathfrak{P}_{2}(S)$, et $S\longrightarrow\mathfrak{P}(A)$,
et cette application est injective sauf si $\Gamma$ est isomorphe au graphe
$\bullet\!\!-\!\!\bullet$ ]. Soient $(U_{s})_{s\in S}$ et
$(E_{a})_{a\in A}$ deux familles d'ensembles, et posons
\[
  X=\coprod_{s\in S}U_{s}\ \amalg\ \coprod_{a\in A}E_{a} .
\]
Pour tout $a_{0}\in A$, considérons les deux comp.\ conn.\ $\Gamma'$,
$\Gamma''$ de l'arbre $\Gamma$ “coupé en $a_{0}$”, et considérons la
décomposition binaire $X=X'\cup X''$ de $X$, où
\[
  X'=\coprod_{s\in S'}U_{s}\amalg\coprod_{a\in A'}E_{a},\qquad
  X''=\coprod_{s\in S''}U_{s}\amalg\coprod_{a\in A''}E_{a},
\]
\note{il n'est pas précisé si $A'$, $A''$ contiennent $a_{0}$ ; la formule
suivante le demande}
donc
\[
  X'\cap X''=E_{a_{0}} .
\]
Cette décomposition est “propre” ssi $X'\neq E_{a_{0}}$,
$X''\neq E_{a_{0}}$. Ceci est le cas, quel que soit le choix de $a_{0}\in
A$, ssi on a la propriété suivante :

($*$) Si $s\in S$ est un sommet-bout, alors $U_{s}\neq\emptyset$

\page{28}
De plus si $a_{1},a_{2}\in A$, $a_{1}\neq a_{2}$, on voit de suite que les
décompositions $\Pi_{a_{1}}$, $\Pi_{a_{2}}$ sont non parallèles, à cela près
qu'il \uncertain{faut exiger} (\uncertain{**}) \uncertain{si $s\in S$}
\struck{\ill{} \ill{} sont \ill{} elles} \ill{} d'ordre 2, alors un des trois
ens.\ $U_{s}$, $E_{a}$, $E_{b}$ \uncertain{est $\neq\emptyset$ si $a$, $b$
sont les deux arêtes en $s$}.
\marginal{\uncertain{si} $a\neq b$ \uncertain{ensemble}, donc
$\Pi_{a}\neq\Pi_{b}$}
\note{la condition (**) est écrite en partie dans l'interligne et en partie
dans la marge gauche, en oblique, où on lit « $E_{a}$, $E_{b}$ \ill{}
$\neq\emptyset$ \ill{} $a$, $b$ sont deux arêtes en $s$ » ; une ligne
biffée court sous l'interligne}

Une façon de visualiser la construction \uncertain{est via} un ensemble
topologique
\[
  \mathcal{X}=(\boldsymbol{\Gamma},S)
\]
(tel que $\Gamma$ en soit une description combinatoire). On fixe un pt sur
chaque arête, de façon à trouver un graphe topologique (subdivision
barycentrique)
\[
  \mathcal{X}^{*}=(\boldsymbol{\Gamma}^{*},\underbrace{S\amalg A'}_{S^{*}}),
\]
où $A'\subset\boldsymbol{\Gamma}$ est en correspondance biunivoque
\uncertain{canonique} avec \struck{\ill{}} $A$.
\note{le $\mathcal{X}$ de la première formule remplace une lettre biffée,
peut-être $\mathcal{A}$. Le $\Gamma$ qu'il souligne (le graphe comme espace
topologique, par opposition à sa description combinatoire $\Gamma$) est
transcrit $\boldsymbol{\Gamma}$ dans toute la page, pour ne pas le confondre avec
le soulignement d'une lecture douteuse}

\struck{\ill{}} [\emph{NB} $A'$ est formé de \struck{pts} sommets d'ordre 2,
et \uncertain{chaque} arête de $\boldsymbol{\Gamma}^{*}$ a une
\uncertain{et une seule} des deux extrémités \uncertain{dans} $A'$,
\uncertain{l'autre} \ill{} $S$ — voici une description complète des
propriétés particulières des $\Gamma^{*}$…]

Alors les points $a'$ de $A'$ coupent $\boldsymbol{\Gamma}$ en deux
composantes connexes, dont les adhérences définissent une décomposition
binaire de $\boldsymbol{\Gamma}$, d'équateur $\{a'\}$. Deux telles
décompositions de $\boldsymbol{\Gamma}$, pour $a'\neq a''$, définissent des
décompositions binaires \emph{strict\add{ement}} non parallèles.

\page{29}
Ici, si on regarde les traces de ces décompositions binaires sur
$S^{*}=S\amalg A'$. Donc pour toute application \emph{surjective}
$X\longrightarrow S^{*}$, si on regarde les images inverses des décompositions
binaires précédentes de $S^{*}$. Or la donnée d'une telle application revient
à celle d'une famille $(X_{\sigma})_{\sigma\in S^{*}}$
\add{d'ensembles} \uncertain{deux à deux disjoints $\neq\emptyset$}, i.e.\ à
celle des deux familles $(U_{s})_{s\in S}$ et $(E_{a})_{a\in A}$ d'ens.\ non
vides.
\note{la première phrase reste sans verbe principal ; « deux à deux
$\neq\emptyset$ » est ajouté dans l'interligne et rattaché par un trait}

Le cas où certains des ensembles $U_{s}$, $E_{a}$ pourraient être vides est à
examiner un peu de plus près, il donne \uncertain{évidemment} \ill{} en
\uncertain{général} \uncertain{encore} ssi on a les conditions ($*$) ($**$)
ci-dessus, et la condition de non-parallélisme \uncertain{strict} ssi on a
\emph{aussi} ($***$) (\uncertain{en plus} que ($**$))

($***$) Si $s\in S$ est un sommet \emph{d'ordre 2} \struck{et}, $a\in A$
une arête incidente, on a \struck{que $U_{s}=\emptyset$, et}
\[
  U_{s}\neq\emptyset\quad\text{ou}\quad E_{a}\neq\emptyset .
\]

\page{31}
\section{Décompositions $n$-aires d'un espace topologique}
\marginal{Version II}
\note{le titre est de sa main, en tête de la page 31 ; « Version II » est
écrit en oblique dans le coin supérieur gauche, et souligné}

1.1. Soit $X$ un espace topologique. On appelle \emph{décomposition binaire}
de $X$ une paire (\uncertain{notation} indicielle\uncertain{ment})
\[
  \Pi=(\Sigma^{\varepsilon})_{\varepsilon\in\omega}\qquad(\omega\text{ ens.\ à deux éléments})
\]
de \struck{ensembles} parties \emph{fermées} (distinctes) de $X$, de réunion
$X$
\[
  X=\bigcup_{\varepsilon\in\omega}\Sigma^{\varepsilon} .
  \tag{1}
\]
On suppose que
\[
  \Sigma^{\varepsilon}\neq\emptyset\qquad(\forall\varepsilon\in\omega)
  \tag{2}
\]
ce qui signifie aussi que
\[
  \Pi\neq\{\emptyset,X\}
  \tag{2'}
\]
et implique que $X\neq\emptyset$, et même $\operatorname{card} X\geqslant2$,
puisque l'on suppose
\[
  \Sigma^{\varepsilon}\neq\Sigma^{\varepsilon'}\quad\text{si}\quad
  \varepsilon,\varepsilon'\in\omega,\ \varepsilon\neq\varepsilon' .
  \tag{3}
\]
On n'exclut pas le cas où on aurait une inclusion
\[
  \Sigma^{\varepsilon'}\subset\Sigma^{\varepsilon}
\]
ce qui équivaut, par (1), à la relation
\[
  \Sigma^{\varepsilon}=X .
\]
Les $\Sigma^{\varepsilon}$ ($\varepsilon\in\omega$) s'appellent les
\emph{hémisphères} de la décomposition binaire, \struck{\ill{}} et l'ensemble
\[
  E=\bigcap_{\varepsilon\in\omega}\Sigma^{\varepsilon}
  \tag{4}
\]
s'appelle l'\emph{équateur} de la décomposition binaire.
\note{dans (4) l'exposant est écrit $\omega$ au lieu de $\varepsilon$}

On n'exclut pas le cas où l'on aurait $E=\emptyset$, ce qui indique que les
$\Sigma^{\varepsilon}$ ($\varepsilon\in\omega$) \uncertain{complémentaires
l'un de l'autre, donc} sont à la fois fermés et ouverts — donc que $X$ est
disconnexe.
\note{« ($\varepsilon\in\omega$) » et « complémentaires l'un de l'autre,
donc » sont ajoutés dans l'interligne, le second d'une lecture incertaine}

\page{32}
Notons les implications
\[
  \left\{
  \begin{array}{l}
    X\ \text{connexe}\ \Longrightarrow\ E\overset{\text{déf}}{=}
    \bigcap_{\varepsilon\in\omega}\Sigma^{\varepsilon}\neq\emptyset\\
    \text{Si les }\Sigma^{\varepsilon}\text{ connexes, alors }
    X\ \text{connexe}\ \Longleftrightarrow\ E\neq\emptyset .
  \end{array}
  \right.
  \tag{5}
\]
Nous poserons
\[
  \Sigma^{0\varepsilon}=\Sigma^{\varepsilon}\smallsetminus E
  =\Sigma^{\varepsilon}\smallsetminus\Sigma^{\varepsilon}\cap
  \Sigma^{\varepsilon'}=X\smallsetminus\Sigma^{\varepsilon'}
  \qquad(\omega=\{\varepsilon,\varepsilon'\})
  \tag{6}
\]
de sorte que l'on a la décomposition
\[
  X=\Sigma^{0\varepsilon_{1}}\cup E\cup\Sigma^{0\varepsilon_{2}}
  \qquad(\omega=\{\varepsilon_{1},\varepsilon_{2}\})
  \tag{7}
\]
de $X$ en \struck{\ill{}} ens.\ mutuellement disjoints.
\note{dans (7) les indices sont d'abord écrits $\varepsilon_{1}$,
$\varepsilon_{2}$, puis corrigés en $\varepsilon'_{1}$, $\varepsilon'_{2}$
dans la parenthèse ; la lecture des indices est incertaine}

Notons
\[
  \overline{\Sigma^{0\varepsilon}}\cap E=\overline{\Sigma^{0\varepsilon}}\cap
  \Sigma^{\varepsilon'}=\overline{\Sigma^{0\varepsilon}}\smallsetminus
  \Sigma^{0\varepsilon}=\dot{\Sigma}^{0,\varepsilon} ,
  \tag{8}
\]
\struck{dans le cas} \struck{on s'intéresse à la cond.}

1.2. Deux décompositions binaires
\[
  \Pi_{1}=\{\Sigma_{1}^{\varepsilon}\}_{\varepsilon\in\omega_{1}},\qquad
  \Pi_{2}=\{\Sigma_{2}^{\varepsilon}\}_{\varepsilon\in\omega_{2}}
  \tag{9}
\]
sont dites \emph{non parallèles}, si elles sont distinctes, et s'il existe
un $\Sigma_{1}^{+}\in\Pi_{1}$, et un $\Sigma_{2}^{+}\in\Pi_{2}$, tels que
l'on ait
\[
  \Sigma_{1}^{+}\cap\Sigma_{2}^{+}=\emptyset
  \tag{10}
\]
condition qui équivaut à :
\[
  \Sigma_{1}^{+}\subset\Sigma_{2}^{-}
  \tag{11}
\]
ou encore :
\[
  \Sigma_{2}^{+}\subset\Sigma_{1}^{-}
  \tag{11 bis}
\]
où $\Sigma_{1}^{-}$, $\Sigma_{2}^{-}$ sont définis par
\[
  \Pi_{1}=\{\Sigma_{1}^{+},\Sigma_{1}^{-}\},\qquad
  \Pi_{2}=\{\Sigma_{2}^{+},\Sigma_{2}^{-}\}
  \tag{12}
\]
\note{le numéro de la dernière formule, mal formé, se lit (12) ou (17) ;
« si elles sont distinctes, et » est ajouté dans l'interligne}

\page{33}
Posons
\[
  \Sigma_{12}=\Sigma_{1}^{-}\cap\Sigma_{2}^{-},\qquad
  E_{1}=\Sigma_{1}^{+}\cap\Sigma_{1}^{-},\qquad
  E_{2}=\Sigma_{2}^{+}\cap\Sigma_{2}^{-}
  \tag{13}
\]
alors on a le diagramme d'inclusions de parties fermées de $X$
\note{c'est la formule (14) de la page ; les flèches sont les inclusions}

\begin{tikzcd}[column sep=small, row sep=small]
 & & X & & \\
 & \Sigma_{2}^{-} \arrow[ur] & & \Sigma_{1}^{-} \arrow[ul] & \\
 \Sigma_{1}^{+} \arrow[ur] & & \Sigma_{12} \arrow[ul] \arrow[ur] & &
 \Sigma_{2}^{+} \arrow[ul] \\
 & E_{1} \arrow[ul] \arrow[ur] & & E_{2} \arrow[ul] \arrow[ur] &
\end{tikzcd}

\note{sous $E_{1}$, une flèche hachurée d'un trait épais, et plus bas une
lettre biffée, peut-être $E$}
\marginal{NB $\Sigma_{1}^{+}\cap\Sigma_{2}^{+}=\emptyset$ donne
$E_{1}\cap\Sigma_{2}^{+}=\emptyset$, $E_{2}\cap\Sigma_{1}^{+}=\emptyset$
i.e.\ $E_{1}\subset\Sigma_{2}^{0-}$, $E_{2}\subset\Sigma_{1}^{0-}$,
$E_{1}\cap E_{2}=\emptyset$}

où les trois carrés sont cartésiens et cocartésiens, i.e.
\[
  \left|
  \begin{array}{l}
    E_{1}=\Sigma_{1}^{+}\cap\Sigma_{12},\quad
    \Sigma_{2}^{-}=\Sigma_{1}^{+}\cup\Sigma_{12}
  \end{array}
  \right.
\]
\note{« cartésiens » est ajouté au-dessus de la ligne, après
« cocartésiens » ; sous la formule, une lettre biffée}
et les deux relations analogues correspondant aux deux autres carrés
\[
  \left\{
  \begin{array}{l}
    E_{2}=\Sigma_{12}\cap\Sigma_{2}^{+},\quad
    \Sigma_{1}^{-}=\Sigma_{12}\cup\Sigma_{2}^{+}\\[4pt]
    \Sigma_{12}\overset{(13)}{=}\Sigma_{1}^{-}\cap\Sigma_{2}^{-},\quad
    X=\Sigma_{2}^{-}\cup\Sigma_{1}^{-} .
  \end{array}
  \right.
  \tag{15}
\]
\note{dans la première ligne de (15), les indices de $\Sigma_{12}$ sont
récrits par-dessus d'autres}

Notons que
\[
  \left\{
  \begin{array}{ll}
    \text{a)} & \Sigma_{1}^{+}\cap\Sigma_{2}^{-}=\Sigma_{1}^{+}\neq\emptyset\\
    \text{b)} & \Sigma_{2}^{+}\cap\Sigma_{1}^{-}=\Sigma_{2}^{+}\neq\emptyset\\
    \text{c)} & \Sigma_{1}^{-}\cap\Sigma_{2}^{-}=\Sigma_{12}\neq\emptyset ,
  \end{array}
  \right.
  \tag{16}
\]
\emph{Prouvons} la dernière relation (les deux premières étant évidentes).
Si on avait $\Sigma_{12}=\emptyset$, comme $E_{1},E_{2}\subset\Sigma_{12}$,
on aurait $E_{1}=\emptyset$, $E_{2}=\emptyset$, de sorte que $\Pi_{1}$ et
$\Pi_{2}$ sont deux \emph{partitions} de $X$ en deux parties non vides, i.e.\
\struck{$\Sigma_{1}^{-}=X\smallsetminus\Sigma_{1}^{+}$} on aurait
\[
  \Sigma_{1}^{-}=X\smallsetminus\Sigma_{1}^{+},\qquad
  \Sigma_{2}^{-}=X\smallsetminus\Sigma_{2}^{+}
\]

\page{34}
d'où
\[
  \Sigma_{12}=X\smallsetminus(\Sigma_{1}^{+}\cup\Sigma_{2}^{+})
\]
donc $\Sigma_{12}=\emptyset$ signifierait que
\[
  X=\Sigma_{1}^{+}\amalg\Sigma_{2}^{+}
\]
et par suite
\[
  \Pi_{1}=\{\Sigma_{1}^{+},\underset{\Sigma_{1}^{-}=X\smallsetminus
  \Sigma_{1}^{+}}{\underset{\shortparallel}{\Sigma_{2}^{+}}}\},\qquad
  \Pi_{2}=\{\Sigma_{2}^{+},\Sigma_{1}^{+}\}
\]
donc \struck{\ill{}} $\Pi_{1}=\Pi_{2}$ contrairement à l'hypothèse.
\note{sous $\Sigma_{2}^{+}$, l'indice de $\Sigma^{-}$ est récrit ; on lit
$\Sigma_{1}^{-}$}

On trouve donc
\marginal{\struck{Corollaire}}

\emph{Proposition} \uncertain{Soient} $\Pi_{1}$, $\Pi_{2}$ deux
décompositions binaires de $X$, mutuellement non parallèles. Alors il existe
un \emph{unique} couple
\[
  (\Sigma_{1}^{+},\Sigma_{2}^{+})\in\Pi_{1}\times\Pi_{2}
\]
tel que l'on ait
\[
  \Sigma_{1}^{+}\cap\Sigma_{2}^{+}=\emptyset .
\]
\note{« Proposition » est écrit au-dessus de « Corollaire », biffé ; l'énoncé
est marqué d'un trait vertical dans la marge}

Soit alors
\[
  C=\Pi_{1}\amalg\Pi_{2}=\omega_{1}\amalg\omega_{2}
  \tag{17}
\]
l'ensemble somme, qui a quatre éléments. On a une application canonique
\[
  C\longrightarrow\mathfrak{P}_{\text{fermé}}(X)
  \tag{18}
\]
qui a comme image
\[
  \{\Sigma_{1}^{+},\Sigma_{1}^{-},\Sigma_{2}^{+},\Sigma_{2}^{-}\},
  \tag{19}
\]
et laquelle on \uncertain{identifie} si (18) est injective i.e.
\note{les numéros (17), (18), (19) sont récrits sur d'autres ; la phrase se
poursuit à la page suivante}

\page{35}
si \struck{\ill{}} cardinal de l'ens.\ (18) est 4, i.e.\ si les quatre
parties de $X$
\[
  \Sigma_{1}^{+},\ \Sigma_{1}^{-},\ \Sigma_{2}^{+},\ \Sigma_{2}^{-}
\]
sont distinctes. A \uncertain{priori} \uncertain{on sait} que l'on sait déjà
par \struck{hypothèse} \struck{($\Pi_{1}$, $\Pi_{2}$ étant} \emph{définition}
des déc.\ binaires
\[
  \Sigma_{1}^{+}\neq\Sigma_{1}^{-},\qquad\Sigma_{2}^{+}\neq\Sigma_{2}^{-}
  \tag{$*$}
\]
et qu'on a
\[
  \Sigma_{1}^{+}\neq\Sigma_{2}^{+}
  \tag{$**$}
\]
puisque $\Sigma_{1}^{+}$ et $\Sigma_{2}^{+}$ sont des parties disjointes
\emph{non vides}. Notons donc qu'on \uncertain{lit} sur le diagramme (14)
\[
  \left.
  \begin{array}{ll}
    \text{a)} & (\Sigma_{1}^{+}=\Sigma_{2}^{-})\Longleftrightarrow
    E_{1}=\Sigma_{12}\qquad(\Rightarrow E_{2}=\emptyset)\\
    \text{b)} & (\Sigma_{2}^{+}=\Sigma_{1}^{-})\Longleftrightarrow
    E_{2}=\Sigma_{12}\qquad(\Rightarrow E_{1}=\emptyset)
  \end{array}
  \right.
  \tag{$*{*}*$}
\]
\note{dans a), le signe entre $\Sigma_{1}^{+}$ et $\Sigma_{2}^{-}$, et entre
$E_{1}$ et $\Sigma_{12}$, est un $=$ repassé en $\neq$ ou l'inverse ; b) et
la suite de l'argument demandent $=$}
et que les deux relations a), b) \uncertain{énergiques} \uncertain{de}
($*{*}*$) ne peuvent être satisfaites simultanément, puisque
\struck{$\Sigma_{12}$} l'on aurait $\Sigma_{12}=E_{1}=E_{2}=\emptyset$
puisque $E_{1}\cap E_{2}=\emptyset$, ce qui contredit (16) (c).

Enfin
\[
  (\Sigma_{1}^{-}=\Sigma_{2}^{-})\Longleftrightarrow
  (\Sigma_{1}^{-}=\Sigma_{2}^{-}=X)\quad\text{i.e.}\quad\Sigma_{12}=X
  \tag{$*{*}{*}*$}
\]
\uncertain{ce} cas est mutuellement exclusif avec chacun des cas ($*{*}*$),
\struck{\ill{} on aurait \ill{}} \struck{$\Sigma_{1}^{-}=\Sigma_{2}^{-}=X$}
$\Sigma_{12}=X$ et $E_{1}=\Sigma_{12}$ \uncertain{on aurait} $E_{1}=X$ donc
$\Sigma_{1}^{+}=\Sigma_{1}^{-}=X$, \struck{\ill{}} ce qui contredit (3).
\note{deux lignes de ce passage sont en partie biffées et récrites dans
l'interligne ; la lecture de leur enchaînement est incertaine}

\emph{Proposition} Pour que l'on ait
$\operatorname{card}\{\Sigma_{1}^{+},\Sigma_{1}^{-},\Sigma_{2}^{+},
\Sigma_{2}^{-}\}=4$, il faut et il suffit que l'on ne soit pas dans l'un des
trois cas suivants, qui s'excluent mutuellement
\[
  \left\{
  \begin{array}{lll}
    \text{a)} & E_{1}=\Sigma_{12} & (\Rightarrow E_{2}=\emptyset)\\
    \text{a')} & E_{2}=\Sigma_{12} & (\Rightarrow E_{1}=\emptyset)\\
    \text{b)} & \Sigma_{12}=X\ \text{i.e.}\ \Sigma_{1}^{-}=\Sigma_{2}^{-}=X &
    (\text{d'où }E_{1}=\Sigma_{1}^{+},\ E_{2}=\Sigma_{2}^{+})
  \end{array}
  \right.
  \tag{20}
\]
\uncertain{auquel cas le cardinal de \ill{} est égal à trois}.
\note{cette dernière ligne est coupée par le bord inférieur de la feuille ;
le numéro (20) est récrit sur un autre}

\page{36}
\emph{Cor}\add{ollaire} \struck{Dans} Dans le cas où
$\operatorname{Card}\{\ \}=4$, et $\Sigma_{1}^{-}\neq X$,
$\Sigma_{2}^{-}\neq X$, les seules relations d'inclusions entre les ensembles
$\Sigma_{1}^{+},\Sigma_{1}^{-},\Sigma_{2}^{+},\Sigma_{2}^{-}$ sont les
relations
\[
  (*)\quad\left\{
  \begin{array}{l}
    \Sigma_{1}^{+}\subset\Sigma_{2}^{-}\\
    \Sigma_{2}^{+}\subset\Sigma_{1}^{-} .
  \end{array}
  \right.
\]
\note{« Cor » est écrit au-dessus d'un « D » biffé ; « et
$\Sigma_{1}^{-}\neq X$, $\Sigma_{2}^{-}\neq X$ » est ajouté au-dessus de la
ligne}
\marginal{Sous forme : Pour que les relations \uncertain{ci-dessus} soient
les seules relations d'inclusion \uncertain{non triviales} entre les
ensembles $\Sigma_{1}^{+},\Sigma_{1}^{-},\Sigma_{2}^{+},\Sigma_{2}^{-}$, il
faut et il suffit que ceux-ci \uncertain{soient} distincts et que l'on ait
$\Sigma_{1}^{-}\neq X$, $\Sigma_{2}^{-}\neq X$.}
\note{note de sa main, en oblique dans la marge gauche, qui reformule le
corollaire ; lecture incertaine par endroits}

En effet, on n'a pas les inclusions inverses, \ill{} par
hypothèse. D'autre part, on n'a \emph{pas}
\[
  \Sigma_{1}^{+}\subset\Sigma_{1}^{-}\quad\text{ou}\quad
  \Sigma_{1}^{-}\subset\Sigma_{1}^{+},
\]
car on aurait \struck{$\Sigma_{1}^{+}$, $\Sigma_{1}^{-}$}
$\Sigma_{1}^{-}=X$, de même on n'a pas
$\Sigma_{2}^{+}\subset\Sigma_{2}^{-}$ ni $\Sigma_{2}^{-}\subset\Sigma_{2}^{+}$,
et on n'a pas $\Sigma_{1}^{+}\subset\Sigma_{2}^{+}$ ni l'inclusion inverse,
puisque $\Sigma_{1}^{+}$ et $\Sigma_{2}^{+}$ sont disjoints et
$\neq\emptyset$. Il reste à voir que l'on n'a pas une inclusion
$\Sigma_{1}^{-}\subset\Sigma_{2}^{-}$ ou l'inclusion inverse, or on
\uncertain{lit} sur le diagramme (14) que \struck{cela} la première
\uncertain{p.\ ex.} \struck{signifierait que
$\Sigma_{12}\simeq\Sigma_{1}^{-}$} \struck{et} \uncertain{implique}
$\Sigma_{2}^{-}=X$, cas que l'on a exclu.
\note{les deux mots entre « la première p.\ ex.\ » et « $\Sigma_{2}^{-}=X$ »
sont récrits sur la ligne biffée ; « la première p.\ ex.\ » est ajouté dans
l'interligne}

\emph{Déf} On dit que $\Pi_{1}$, $\Pi_{2}$ sont \emph{strictement non
parallèles} si elles sont non parallèles, et si les quatre ensembles
$\Sigma_{1}^{+},\Sigma_{1}^{-},\Sigma_{2}^{+},\Sigma_{2}^{-}$ sont
distincts, et si de plus $\Sigma_{1}^{-}\neq X$, $\Sigma_{2}^{-}\neq X$ —
i.e.\ si les seules relations d'inclusion entre ces ensembles
$\Sigma_{1}^{+},\Sigma_{1}^{-},\Sigma_{2}^{+},\Sigma_{2}^{-}$ sont
$\Sigma_{1}^{+}\subset\Sigma_{2}^{-}$, $\Sigma_{2}^{+}\subset\Sigma_{1}^{-}$ —
ce qui signifie aussi que l'on a :
\[
  \Sigma_{1}^{-}\neq X,\ \Sigma_{2}^{-}\neq X\ \ (\text{on dit donc que
  }\Pi_{1},\Pi_{2}\text{ sont \emph{propres}}),\quad
  E_{1}\neq\Sigma_{12},\ E_{2}\neq\Sigma_{12} .
  \tag{21}
\]
\note{la définition est marquée d'un trait vertical dans la marge ; « si
elles sont non parallèles, et » est ajouté dans l'interligne}

\page{37}
Mais si cette condition n'est pas satisfaite pour $\Pi_{1}$, $\Pi_{2}$ non
parallèles, on \uncertain{écrit} \ill{}, et on regarde l'ens.\ à quatre
éléments
\[
  C=\Pi_{1}\amalg\Pi_{2}=\omega_{1}\amalg\omega_{2}=\{\varepsilon_{1}^{+},
  \varepsilon_{1}^{-},\varepsilon_{2}^{+},\varepsilon_{2}^{-}\}
  \tag{17}
\]
et la relation d'ordre dessus, où les seules relations $<$ sont données
par
\[
  \varepsilon_{1}^{+}<\varepsilon_{2}^{-},\qquad
  \varepsilon_{2}^{+}<\varepsilon_{1}^{-}
  \tag{22}
\]

\begin{tikzcd}[column sep=small, row sep=small]
\varepsilon_{1}^{-} & \varepsilon_{2}^{-} \\
\varepsilon_{1}^{+} \arrow[ur] & \varepsilon_{2}^{+} \arrow[ul]
\end{tikzcd}

\note{les quatre éléments sont des points ; sous $\varepsilon_{1}^{+}$ un
mot biffé}

1.3. Soit \struck{$(\Pi_{\alpha})_{\alpha\in A}$ une famille}
\[
  (\Pi_{\alpha})_{\alpha\in A},\qquad
  \Pi_{\alpha}=\{\Sigma_{\alpha}^{\varepsilon}\}_{\varepsilon\in\omega_{\alpha}}
  \tag{23}
\]
\struck{une famille} \add{un} ens.\ (\uncertain{notation}
indicielle\uncertain{ment}) \struck{de} \uncertain{de} décompositions
binaires de $X$, \uncertain{on suppose} \struck{parallèles}
\[
  \forall\alpha,\beta\in A,\ \alpha\neq\beta,\quad\Pi_{\alpha},\Pi_{\beta}
  \text{ sont mut\add{uellemen}t non parallèles.}
  \tag{24}
\]
\emph{Déf} \struck{On} On dit alors que $(\Pi_{\alpha})_{\alpha\in A}$ est
non parallèle. (Terminologie analogue pour \emph{strictement} non
parallèle).
\marginal{c'est un élément de
$\mathfrak{P}(\mathfrak{P}_{2}(\mathfrak{P}(X)))$ \ill{}}
\marginal{On dit, si $\operatorname{card} A=n$, que l'on a une
\emph{décomposition $n$-aire} non parallèle de $X$}
\note{deux notes de sa main, en oblique dans la marge gauche ; la seconde
porte le titre de la page 31}

Considérons
\[
  C=\coprod_{\alpha\in A}\omega_{\alpha}
  \tag{25}
\]
avec l'involution sans pt fixe $\sigma$ qui sur chaque $\omega_{\alpha}$
induit une transposition, et sur $C$ la relation d'ordre, dans laquelle on
a
\note{la page s'arrête sur ce « a », en bas de feuille ; l'ajout sur
l'involution $\sigma$ est écrit dans l'interligne à droite de (25)}

\page{38}
\[
  \Bigl(x<y,\ x\in\omega_{\alpha},\ y\in\omega_{\beta}\Bigr)
  \Longleftrightarrow
  \Bigl(\alpha\neq\beta,\ x\underset{\{\Pi_{\alpha},\Pi_{\beta}\}}{<}y\Bigr)
  \tag{26}
\]
\note{dans le membre de gauche, « $x\in\omega_{\alpha}$ »,
« $y\in\omega_{\beta}$ » sont écrits verticalement, $\omega_{\alpha}$ sous
$x$ et $\omega_{\beta}$ sous $y$, avec un $\in$ couché ; sous l'accolade du
membre de droite : « i.e.\
$\Sigma_{\alpha}^{x}\cap\Sigma_{\beta}^{\sigma y}=\emptyset$ »}

\marginal{il faut vérifier que c'est \uncertain{transitif} \ill{} relation
\ill{} \uncertain{i.e.} \uncertain{d'ordre}}
\note{note en oblique dans la marge gauche, sur le mot « Alors »}

Alors $\sigma$ est un antiautomorphisme (involutif sans pt fixe) de $C$.
L'application
\[
  \begin{array}{c}
    C\overset{\varphi}{\longrightarrow}\mathfrak{P}_{\text{fermé}}(X)\\
    (\alpha,x)\longmapsto\Sigma_{\alpha}^{x}
  \end{array}
  \tag{27}
\]
est une application croissante. Pour qu'elle soit injective, et définisse un
isomorphisme d'ens.\ ordonnés de $C$ sur son image, il faut et il suffit que
$(\Pi_{\alpha})_{\alpha\in A}$ soit \emph{strictement non parallèle}.

Inversement on a :

\emph{Définition} On appelle \emph{quartette \struck{ordonné}}
\add{\uncertain{ordonné}} (de décompositions $n$-aires) un ens.\
\uncertain{fini} \struck{\ill{}} $C$, muni d'une involution $\sigma$ sans pt
fixe, et d'une relation d'ordre, ayant les propriétés suivantes :
\note{« quartette » est souligné ; « ordonné » y est biffé puis, semble-t-il,
récrit ; au-dessus de « un ens.\ » un mot ajouté, peut-être « fini »}

\struck{a) Si $x,y\in C$ et $x<y$, alors $\sigma x$ et $\sigma y$ …}
\struck{Si $x,y\in C$, $\sigma x\neq y$, $x$ et $y$ …}
\struck{on n'a pas $x<y$ \ill{}}
\note{trois lignes biffées d'un trait chacune, et le numéro a) barré d'une
grande croix}

a) si $x,y\in C$, \struck{\ill{}} alors
\uncertain{$\bigl\{$}$x,y$ comparables
\struck{$\sigma x,y$ et $x,\sigma y$} ssi $\sigma x,y$ non comparables,
\struck{et $\sigma y<\sigma x$}

\struck{b) si $x,y\in C$, alors $x<y\Rightarrow\sigma y<\sigma x$}
\note{b) est biffé d'un trait, mais son contenu est repris dans le NB qui
suit ; la numérotation a), b) est corrigée plusieurs fois dans la marge}

\marginal{NB \ill{} que $x<y$ \uncertain{ssi} \ill{} $x<y$ \uncertain{non}
\uncertain{comparables} \ill{} $\sigma$}
\note{note en oblique dans la marge gauche, en bas de page ; lecture très
incertaine}

\emph{NB} \uncertain{Cela} \ill{} \uncertain{que} $\sigma$ est un
antiautomorphisme de l'ens.\ ordonné $C$
\note{la dernière ligne de la page s'arrête là}

\page{39}
\emph{Proposition} Soit $(C,\sigma,<)$ un \struck{\ill{}}
quartette ordonné. Alors pour tout $x\in C$, \struck{les ens.}
\[
  C_{\geqslant x}=\{y\in C\mid y\geqslant x\},\qquad
  C_{\leqslant x}=\{y\in C\mid y\leqslant x\}
\]
sont \emph{totalement} ordonnés.
\note{l'énoncé est marqué d'un trait vertical dans la marge ; le mot
souligné est abrégé « tot$^{t}$ »}

Il suffit de \uncertain{prouver} la première (par dualité), \struck{soit
donc} ce qui signifie aussi :
\[
  \text{si } x<y,\ x<z,\ \text{avec } y\neq z,\ \text{alors } (y,z)\ \text{comparables.}
\]
\note{à gauche, un petit croquis : trois points $x$ (en bas), $y$ et $z$
(en haut), avec des flèches $x\to y$, $x\to z$ et deux traits croisés}

[On note que $y$ et $z$ sont non conjugués \uncertain{mod} $\sigma$, car
\uncertain{ce serait} \uncertain{contre} (28 a),] Si $y$ et $z$ étaient
\emph{in}comparables, on aurait que $\sigma y$ et $z$ sont comparables,
\uncertain{or} on ne peut avoir $z\leqslant\sigma y$ car on aurait
$x\leqslant\sigma y$, ce qui contredit $x<y$ par (28 a). On
\struck{\ill{}} aurait donc $\sigma y<z$ d'où $\sigma z<$
\note{la phrase s'arrête là ; le reste de la page est blanc. La référence
(28 a) renvoie sans doute à la condition a) de la page 38, dont le numéro
n'est pas lisible sur la page ; « $y$ et $z$ » après « Si » sont récrits
sur d'autres lettres}

\end{document}
